На 48 независимо сгенерированных обучающих наборах ни один эпизод обучения или валидации не достигает заранее заданного критерия тяжёлого кризиса. В основном сценарии, где нагрузка повышает вероятность повторной работы, подготовка снижает стоимость на 38.2\% относительно планирования по модели без влияния нагрузки на длительность и повторную работу и на 38.1\% относительно добавления исполнителя при 20 незавершённых задачах. Оба парных доверительных интервала с поправкой на множественные сравнения исключают ноль. Превосходство над сокращением повторной работы при том же пороге 20 незавершённых задач не установлено. Описательные нормированные сравнения по четырём механизмам показывают выигрыш выбора действия по модели относительно каждого фиксированного действия, а также проигрыш без обратной связи и чувствительность к стоимости действия.
